% =====================================================================================
% EXPERIMENTAL SPECIFICATION -- FOR REVIEW BEFORE THE v3 PRODUCTION RUN
%
% This appendix states the setup in full so it can be checked before any of it is executed.
% Numbers already measured are marked [measured]; numbers awaiting the production run are
% marked [pending] and are the only thing the run adds.
% =====================================================================================
\section{Experimental specification}\label{app:spec}

\subsection{The two admissibility rules}\label{app:rules}

Everything in this study is governed by two rules. They are stated first because they decide
which methods may appear at all, and because one of them has a clause that is easy to get wrong.

\begin{description}
\item[Rule 1 (hard labels).] A training record emitted after a wait $\Delta$ carries the label
  \begin{equation}\label{eq:rule1}
  y_i \;=\; \ind[\,\text{conv\_ts}_i \;\le\; \text{click\_ts}_i + \Delta\,],
  \end{equation}
  and that value is frozen. A conversion arriving after $\text{click\_ts}_i + \Delta$ must not
  flip it, must not add a second supervision signal, and must not be recovered later.

\item[Rule 1b (teacher soft labels).] A soft target is exempt from~\eqref{eq:rule1} provided the
  model producing it is itself trained under Rule 1. Concretely: the teacher trains only on
  \emph{matured} records, whose hard labels obey~\eqref{eq:rule1} at $\Delta = v$; it is then
  \emph{evaluated} on the \emph{freshest} clicks, those released at $\Delta = o$, and the student
  trains on that read-out. This is admissible because the student's record carries no hard label
  drawn from its own future --- only a deterministic function of its features, computed by
  parameters frozen at emission.

\item[Rule 2 (single visit).] Each click contributes \emph{exactly one} training record to the
  served model, ever.
\end{description}

\paragraph{The causality obligation Rule 1b creates.} Rule 1b is only sound if the teacher is
causal at the moment it scores. At cutoff $T_n$ our teacher has trained on clicks satisfying
$C_i + v \le T_n$ and on nothing else, so every parameter it holds derives from outcomes that were
knowable at $T_n$. Scoring a click released at $T_n$ therefore uses no future information. This is
not assumed: it is enforced in code and checked by gate~G1 of
\texttt{tr\_verify\_distill.py}, which fails the run if any teacher update reads a row whose
window closes after the current cutoff.

\paragraph{What Rule 1b does \emph{not} permit.} The teacher may not be trained on the fresh
stream, and the student may not consume the teacher's read-out for a click whose window has since
closed. The non-compliant reference \textsc{Wait}-hybrid, which trains the teacher on matured
\emph{and} fresh records, measures $2.08$ visits per click and is reported only to bound what the
rules cost.

\subsection{Per-method compliance}\label{app:compliance}

Table~\ref{tab:spec-compliance} classifies every method against both rules. Record counts are
\emph{measured} over a full 720-hour replay, not asserted from the description.

\begin{table}[h]
\caption{Every method against both rules. \textsuperscript{1b} marks a soft target admissible
under Rule~1b: the producing model is itself trained under Rule~1. Records per click are measured
over a full 720-hour replay}
\label{tab:spec-compliance}
\centering\small
\begin{tabular}{l l c@{\hskip 1.1em} c@{\hskip 1.1em} r@{\hskip 1.1em} c}
\toprule
method & target $\tilde y$ & Rule 1a & Rule 2 & rec./click & in run\\
\midrule
\textsc{Fresh} & $\tilde y=\Yo$ & \checkmark & \checkmark & 1.00 & \checkmark\\
\textsc{Reweight} & $\tilde y=\Yo$, loss $\times\,w_{\text{iw}}(x)$ & \checkmark & \checkmark & 1.00 & \checkmark\\
\textsc{Wait} & $\tilde y=\Yv$ at $\Delta{=}v$ & \checkmark & \checkmark & 1.00 & \checkmark\\
\textsc{Twice} & $\tilde y=\Yo$ + delay head & \checkmark & \checkmark & 1.00 & \checkmark\\
\textsc{Correct} & $\tilde y=\Yo+(1-\Yo)\,w(x)$ & \checkmark & \checkmark & 1.00 & \checkmark\\
\textsc{Correct-Distill} & $\tilde y=\Yo+(1-\Yo)\,w^{T}(x)$ & \checkmark\,\textsuperscript{1b} & \checkmark & 1.00 & \checkmark\\
\textsc{Distill-Only} & $\tilde y=\hat p_v(x)$ & \checkmark\,\textsuperscript{1b} & \checkmark & 1.00 & \checkmark\\
\textsc{Oracle} & $\tilde y=\Yv$ at $\Delta{=}o$ & $\times$ & \checkmark & 1.00 & \checkmark\\
\midrule
\textsc{Vanilla} & $\Yo$ plus a late positive & \checkmark & $\times$ & 1.08 & $\times$\\
\textsc{Es-Dfm} & duplicated delayed positive & \checkmark & $\times$ & 1.08 & $\times$\\
\textsc{Defuse} & duplicated stream + correction & \checkmark & $\times$ & 1.08 & $\times$\\
\textsc{Ddfm} & both streams together & \checkmark & $\times$ & 1.89 & $\times$\\
\textsc{Miss} & delayed positive into every head & \checkmark & $\times$ & 2.05 & $\times$\\
\textsc{Ftp} & $K$ maturity-gated tasks + prophet & \checkmark & $\times$ & 3.00 & $\times$\\
\textsc{Pi} & two fixed-window predictors & \checkmark & $\times$ & 2.00 & $\times$\\
\textsc{If-Dfm} & influence update when a label reverses & $\times$ & $\times$ & 2.90 & $\times$\\
\bottomrule
\end{tabular}

\end{table}

\subsection{Common substrate}\label{app:substrate}

\paragraph{Corpus and split.} Criteo-DFD~\cite{chapelle2014modeling}: $15{,}609{,}913$ clicks with
conversion timestamps, of which $3{,}538{,}028$ ever convert at any horizon and the longest
observed delay is $30.1$ days. Pre-train on days $[0,30)$, then stream $[30,60)$ with $719$ hourly
cut-offs. At each cut-off a model consumes the records released in that hour and then predicts
every click of the next hour exactly once, scored retrospectively against $\Yv$.

\paragraph{Seeds do not change the data.} The split, the cut-offs and the per-hour cohorts are
deterministic; a seed varies only parameter initialisation and minibatch order. This was verified
rather than assumed: across all five seeds at a fixed rung the evaluation set is identical
($n_{\text{eval}} = 8{,}196{,}131$, base rate $0.136139$, $719$ cohorts, all bit-identical). There
is therefore one data specification, not five.

\paragraph{Label statistics.} Table~\ref{tab:spec-labels} gives the label composition at both
windows. The arithmetic is asserted in code: positives plus negatives equal $N$, the fresh
positives are a strict subset of the served positives, and the observed rate equals
$\text{base} \times F(o)$ to machine precision.

\begin{table}[h]
\caption{Label composition of the replay at both windows. $\Yo$ does not depend on $v$, so the
fresh label is identical at both; only the target it must predict moves}
\label{tab:spec-labels}
\centering\small
\begin{tabular}{lrr}
\toprule
quantity & $v=1$\,d & $v=7$\,d\\
\midrule
\multicolumn{3}{l}{\emph{clicks in the replay $[0,60)$\,d}}\\
\quad clicks $N$ & 15,609,913 & 15,609,913\\
\quad ever converts (any horizon) & 3,538,028 & 3,538,028\\
\midrule\multicolumn{3}{l}{\emph{served target $\Yv=\ind[D\le v]$}}\\
\quad positives & 2,145,812 & 2,875,121\\
\quad negatives & 13,464,101 & 12,734,792\\
\quad positive rate & 0.1375 & 0.1842\\
\midrule\multicolumn{3}{l}{\emph{fresh label $\Yo=\ind[D\le o]$, $o=1$\,h}}\\
\quad positives & 1,649,879 & 1,649,879\\
\quad negatives & 13,960,034 & 13,960,034\\
\quad positive rate & 0.1057 & 0.1057\\
\quad censored positives ($\Yv{=}1,\Yo{=}0$) & 495,933 & 1,225,242\\
\quad $F(o)=P(D\le o\mid D\le v)$ & 0.7689 & 0.5738\\
\midrule\multicolumn{3}{l}{\emph{delay of within-window converters (hours)}}\\
\quad median & 0.17 & 0.44\\
\quad p90 & 8.56 & 89.67\\
\quad p99 & 23.17 & 160.67\\
\bottomrule
\end{tabular}

\end{table}

Two facts drive much of the paper. First, $\Yo$ does not depend on $v$, so the fresh label is
\emph{identical} at both windows ($1{,}649{,}879$ positives); only the target it is asked to
predict moves. Second, censoring deepens with the window --- $F(o)$ falls from $0.769$ at $v=1$\,d
to $0.574$ at $v=7$\,d --- so the fresh label is a good proxy at one day and a poor one at seven.

\paragraph{Record timing.} Figure~\ref{fig:spec-timeline} shows, for one click, when its features
are known, when each candidate label becomes knowable, and when it is scored.

% Record timing for one click. Greyscale, LNCS text width.
% Laid out so nothing overlaps: the two pipeline rows sit above the axis at distinct heights with
% their markers on the right, the teacher arrow runs BELOW them, and the two early tick labels are
% staggered vertically because C_i and C_i+o are close together on any linear time axis.
\begin{figure}[t]
\centering
\begin{tikzpicture}[font=\scriptsize, x=1cm, y=1cm,
  dot/.style={circle, fill, inner sep=1.15pt},
  ar/.style={draw, -{Latex[length=1.2mm]}, line width=.45pt}]

  % ---- pipeline rows (above the axis) -------------------------------------------------
  % the row text states its own emission time, so no marker is needed on the row -- an inline
  % dot at C_i+o would land inside the label
  \node[anchor=west] at (0,2.30) {$\mathcal{D}_{\text{mat}}$: one record at $C_i{+}v$, label $\Yv_i$, frozen};
  \node[anchor=west] at (0,1.72) {$\mathcal{D}_{\text{fresh}}$: one record at $C_i{+}o$, label $\Yo_i$, frozen};

  % ---- Rule 1b: teacher trained on the matured pipeline, SCORED at the fresh release ---
  \draw[ar] (9.05,1.15) -- node[above, pos=.52]
      {teacher trained on $\mathcal{D}_{\text{mat}}$, \emph{scored} at $C_i{+}o$ (Rule 1b)} (1.62,1.15);

  % ---- time axis ----------------------------------------------------------------------
  \draw[ar] (-0.25,0) -- (10.3,0);
  \node[anchor=south east] at (10.3,.04) {time};
  \foreach \x in {0, 1.55, 4.4, 9.05} { \draw[line width=.45pt] (\x,.11) -- (\x,-.11); }
  % staggered: C_i and C_i+o are adjacent, so their captions sit at different depths
  \node[anchor=north, align=center] at (0,-.14)   {$C_i$\\features};
  \node[anchor=north, align=center] at (1.75,-.70) {$C_i{+}o$\\fresh record};
  \node[anchor=north, align=center] at (4.4,-.14) {$V_i$\\conversion};
  \node[anchor=north, align=center] at (9.05,-.14) {$C_i{+}v$\\matured record};

  % ---- evaluation ----------------------------------------------------------------------
  \draw[dotted, line width=.4pt] (1.55,-1.38) -- (3.15,-1.38);
  \node[anchor=west] at (3.22,-1.38)
      {scored once in $(T_n,T_{n+1}]$, retrospectively against $\Yv$};
\end{tikzpicture}
\caption{Timing for a single click. Features are known at $C_i$. The fresh pipeline writes one
record at $C_i{+}o$ carrying $\Yo$; the matured pipeline writes one at $C_i{+}v$ carrying $\Yv$.
Rule~2 permits the served model only one of them. Under Rule~1b the teacher learns from the
matured pipeline but is \emph{evaluated} at the fresh release, so the student's record carries a
soft target and no future hard label}
\label{fig:spec-timeline}
\end{figure}


\subsection{The A0--A4 capacity ladder}\label{app:ladder}

The previous S0--S4 ladder varied only hash rows against a \emph{fixed} trunk of $143{,}293$
parameters, so total capacity spanned just $66\times$ and was trunk-dominated below the top rung;
it was a table-resolution axis described as a capacity axis. A0--A4 pairs a rung of the
compute-efficient frontier with a row count and lets the embedding width rise monotonically,
giving $\sim\!10\times$ per step over $10{,}346\times$. A4 is byte-identical to S4, so the two
ladders share a top rung.

\begin{table}[h]
\caption{The A0--A4 ladder, introspected from the instantiated models. The trunk grows with the
tables, and embedding width rises monotonically $2\to4\to8$}
\label{tab:spec-ladder}
\centering\small
\begin{tabular}{lrrrrrr}
\toprule
rung & frontier node & emb.\ dim & table rows & tables & trunk & total\\
\midrule
A0 & n0 & 2 & 36 & 648 & 278 & 926\\
A1 & n2 & 4 & 256 & 9{,}216 & 594 & 9{,}810\\
A2 & n4 & 8 & 1{,}280 & 92{,}160 & 2{,}799 & 94{,}959\\
A3 & n6 & 8 & 16{,}384 & 1{,}179{,}648 & 4{,}791 & 1{,}184{,}439\\
A4 & -- & 8 & 131{,}072 & 9{,}437{,}184 & 143{,}293 & 9{,}580{,}477\\
\bottomrule
\end{tabular}

\end{table}

\paragraph{Why the width rises.} Two earlier designs were rejected on measurement. Letting each
frontier node keep its own width put A3 at $\text{emb}=48$ with $2048$ rows, which scored
\emph{worse} than A2 at $\text{emb}=8$ with $1280$ rows despite ten times the parameters: on this
corpus resolution carries the signal and width does not. Pinning the width at $8$ everywhere fixed
that but removed the Oracle inversion entirely ($+0.0166 \to -0.0004$ at the bottom rung), because
the inversion needs genuine representational starvation and $\text{emb}=8$ is already wide enough
at $36$ rows. The adopted ladder raises width monotonically, $2\to4\to8$, keeping the starved
bottom without the mid-ladder misallocation.

\subsection{Experiment matrix}\label{app:matrix}

\begin{table}[h]
\caption{What the production run executes. Seeds replicate initialisation and shuffling only;
the data is identical across them}
\label{tab:spec-matrix}
\centering\small
\begin{tabular}{l l r}
\toprule
axis & levels & count\\
\midrule
method & the 12 arms of Table~\ref{tab:spec-compliance} & 12\\
capacity & A0, A1, A2, A3, A4 & 5\\
window & $v=1$\,d, $v=7$\,d & 2\\
seed & $0\ldots4$ (init and shuffle only; data identical) & 5\\
\midrule
\multicolumn{2}{l}{streaming runs} & 600\\
\multicolumn{2}{l}{learning-rate selections (7 rates{,} per arm/rung/window)} & 840\\
\multicolumn{2}{l}{NE curves recorded (seed 0)} & 120\\
\bottomrule
\end{tabular}

\end{table}

Learning rates are selected \emph{per arm, per rung, per window} from
$\{10^{-6}, 2.5{\times}10^{-6}, 10^{-5}, 5{\times}10^{-5}, 2{\times}10^{-4}, 5{\times}10^{-4},
10^{-3}\}$ on a held-out chronological slice of the pre-training period, and auxiliary learning
rates are swept per method for the reason established in Sect.~\ref{sec:tuning}. Curves are
recorded at seed~0; all reported statistics use five seeds and paired tests over matched seeds.
